%=======================================================================
%  Paper 2 (AI) — Where should learning enter a physics-based dehazing
%  pipeline? A stage-wise study of the Dark Channel Prior.
%
%  Every number comes from the Kaggle run of NB 4 (dehaze-hybrid-stages):
%    numbers_hybrid.tex   <- macros (hyb...)
%    numbers.tex          <- paper 1's macros (NB 2), reused for the scaffold
%    figures/*.pdf, tab_*.png
%  Macros fall back to "??" so the document always compiles.
%  %% AUTHORS sections are to be written by the group.
%=======================================================================
\documentclass[conference]{IEEEtran}
\IEEEoverridecommandlockouts
\usepackage{cite}
\usepackage{amsmath,amssymb}
\usepackage{graphicx}
\usepackage{booktabs}
\usepackage{array}
\usepackage{xcolor}
\usepackage[hidelinks]{hyperref}
\graphicspath{{figures/}}

\InputIfFileExists{numbers.tex}{}{}
\InputIfFileExists{numbers_hybrid.tex}{}{}
\makeatletter
\newcommand{\fb}[1]{\@ifundefined{#1}{\expandafter\newcommand\csname #1\endcsname{??}}{}}
\makeatother
% paper-1 macros used here
\fb{nImages}\fb{nTune}\fb{nTest}\fb{nDatasets}\fb{dataSource}\fb{splitFingerprint}\fb{tunedCfg}\fb{tunedPSNR}\fb{basePSNR}\fb{inputPSNR}
% NB 4 macros
\fb{hybNTest}\fb{hybSeeds}\fb{hybSteps}\fb{hybTrainSource}\fb{hybNRealHazy}\fb{hybDevice}\fb{hybBestStage}\fb{hybBestStageGain}\fb{hybBrightN}
\makeatletter
\def\fbarm#1{\fb{hyb#1PSNR}\fb{hyb#1SSIM}\fb{hyb#1CIEDE}\fb{hyb#1Delta}\fb{hyb#1Lo}\fb{hyb#1Hi}\fb{hyb#1P}\fb{hyb#1Ms}\fb{hyb#1Params}%
  \fb{hyb#1Light}\fb{hyb#1Medium}\fb{hyb#1Dense}\fb{hyb#1Synth}\fb{hyb#1Real}\fb{hyb#1Drop}%
  \fb{hyb#1Bright}\fb{hyb#1Else}\fb{hyb#1Ratio}\fb{hyb#1BDelta}\fb{hyb#1BLo}\fb{hyb#1BHi}\fb{hyb#1BP}}
\makeatother
\fbarm{Input}\fbarm{Clahe}\fbarm{Base}\fbarm{Tuned}\fbarm{HA}\fbarm{HTsup}\fbarm{HTprior}\fbarm{HJ}\fbarm{HAll}\fbarm{Aod}

\setlength{\textfloatsep}{6pt plus 2pt minus 2pt}
\setlength{\floatsep}{6pt plus 2pt minus 2pt}
\setlength{\intextsep}{6pt plus 2pt minus 2pt}
\setlength{\dbltextfloatsep}{6pt plus 2pt minus 2pt}
\setlength{\abovecaptionskip}{3pt}
\setlength{\belowcaptionskip}{0pt}

\begin{document}

\title{Where Should Learning Enter a Physics-Based Dehazing Pipeline?\\
A Stage-Wise Study of the Dark Channel Prior}

\author{
\IEEEauthorblockN{Md.\ Asif Hossain \quad Nabil Subhan \quad K M Nudar}
\IEEEauthorblockA{Department of Computer Science and Engineering, East West University, Dhaka, Bangladesh\\
\{2022-3-60-007, 2022-3-60-063, 2022-3-60-234\}@std.ewubd.edu}
}
\maketitle

\begin{abstract}
Learned dehazers replace the whole physics-based pipeline of the dark channel prior (DCP) with a network, so it is not
known \emph{which} of its stages actually benefits from learning. We keep the DCP pipeline as a fixed classical scaffold and
replace one stage at a time---atmospheric light, transmission refinement, or radiance recovery---with a small residual
module (17--23\,k parameters) trained with an identical budget of \hybSteps{} steps, then compare each hybrid, their
composition and an end-to-end network (AOD-Net) trained on the same data, on the identical test split of
\hybNTest{} images from \nDatasets{} benchmarks used in our companion classical study. All comparisons are paired against the
tuned classical pipeline with bootstrap confidence intervals over \hybSeeds{} seeds. The stage that benefits most from
learning is \hybBestStage{} (\hybBestStageGain{}\,dB); composing all three modules gives \hybHAllDelta{}\,dB and the
end-to-end network \hybAodDelta{}\,dB relative to the tuned DCP. Training the transmission refiner with the dark
channel prior itself as its loss, without ground truth, changes the bright-region error by \hybHTpriorBDelta{}
(95\,\% CI [\hybHTpriorBLo, \hybHTpriorBHi]) whereas the same network trained with ground truth changes it by
\hybHTsupBDelta{}, showing that a prior used as a loss inherits the prior's failure. Gains measured on synthetic haze
transfer to real haze with a drop of \hybHAllDrop{}\,dB for the hybrid versus \hybAodDrop{}\,dB for the end-to-end
model.
\end{abstract}

\begin{IEEEkeywords}
image dehazing, dark channel prior, physics-informed learning, hybrid pipelines, unsupervised loss, ablation study
\end{IEEEkeywords}

%=========================== I ==========================================
\section{Introduction}
Single-image dehazing inverts the atmospheric scattering model $I = J\,t + A(1-t)$~\cite{narasimhan2002}. The dark
channel prior (DCP)~\cite{he2009,he2011} solves it with a fixed sequence of classical stages---dark channel, atmospheric
light $A$, raw transmission, transmission refinement, radiance recovery---while learned methods either predict a stage
from data (DehazeNet~\cite{cai2016} and MSCNN~\cite{ren2016} regress $t$) or replace the whole pipeline
(AOD-Net~\cite{li2017}, FFA-Net~\cite{qin2020}, DehazeFormer~\cite{song2023}). Comparisons between the two families
report end results only, so three questions remain open: which classical stage is actually the bottleneck that
learning should address; whether a prior can serve as a training signal without inheriting its known failure in bright
regions; and whether keeping the physics scaffold helps a model trained on synthetic haze transfer to real haze.

We answer them with a controlled design. The DCP pipeline, tuned in our companion study on a scene-grouped tuning split,
is kept fixed. Three residual modules---$A = A_{\text{cls}} + \Delta A$, $t = t_{\text{guided}} + \Delta t$,
$J = J_{\text{DCP}} + \Delta J$---each replace one stage; an untrained module reproduces the classical pipeline, so any
gain is attributable to learning at that stage. All modules, their composition and an end-to-end AOD-Net are trained
with the same data and step budget and evaluated under the same paired protocol. Our contributions:
\begin{enumerate}
  \item A stage-wise hybridisation of DCP with differentiable classical stages, so each learned module is trained
        \emph{through} the fixed remainder of the pipeline and is an exact drop-in at inference.
  \item A ranking of the stages by learning gain, with paired CIs over seeds, density strata and cost (H4).
  \item A controlled test of the dark channel prior as an unsupervised loss: the same refiner trained with and
        without ground truth, measured inside intrinsically bright regions (H5).
  \item A sim-to-real transfer comparison between single-stage hybrids and an end-to-end model (H6).
\end{enumerate}

%=========================== II =========================================
\section{Related Work}
\textbf{Learning one stage.} DehazeNet~\cite{cai2016} and MSCNN~\cite{ren2016} learn the transmission map and keep the
classical $A$ and recovery; their networks are trained on synthetic transmission targets, not through the pipeline.
\textbf{Learning everything.} AOD-Net~\cite{li2017} folds $t$ and $A$ into one variable; later
architectures~\cite{qin2020,song2023} scale capacity. \textbf{Priors as losses.} Golts \emph{et al.}~\cite{golts2020}
trained a dehazing network with the DCP matting energy as an unsupervised loss; we adopt the simpler
dark-channel-of-output loss and, unlike prior work, test it specifically where the prior is violated. \textbf{Guided
filtering.} The guided filter~\cite{he2010} is differentiable~\cite{wu2018} and is used here as the classical stage the
learned refiner corrects. \textbf{Benchmarks.} RESIDE~\cite{li2019} (ITS for training, SOTS for testing), O-HAZE,
I-HAZE, Dense-Haze and NH-HAZE~\cite{ancuti2018o,ancuti2018i,ancuti2019,ancuti2020} for real haze.

%=========================== III ========================================
\section{Method}
\subsection{Classical scaffold}
The scaffold is the DCP pipeline with the configuration selected on the tuning split in our companion study
(\emph{\tunedCfg}). Its stages are re-implemented differentiably for training: the dark channel is a channel-minimum
followed by min-pooling (grey erosion), the guided filter~\cite{he2010} is four box filters (average pooling), and
recovery is the closed-form $J = (I - A)/\max(t, t_0) + A$. The classical $A$-estimator is called on each training
crop so that the same estimator is used at training and test time. At inference the NumPy stages of the companion
study run and only the learned module executes in PyTorch.

\subsection{Residual modules}
\textbf{H-A} (\hybHAParams{} parameters): a four-stride convolutional encoder on a $128^2$ downsample of the hazy image,
concatenated with the broadcast classical $A$, outputs $\Delta A\in[-0.3,0.3]^3$.
\textbf{H-t} (\hybHTsupParams{}): a five-layer dilated convolutional network (dilations 1--8) on
$[I,\,\tilde t,\,t_{\text{guided}},\,I^{\text{dark}}]$ outputs $\Delta t\in[-0.5,0.5]$, added to the guided-filter
estimate.
\textbf{H-J} (\hybHJParams{}): a four-layer network on $[I,\,J_{\text{DCP}},\,t]$ outputs $\Delta J\in[-0.25,0.25]^3$.
All bounds are $\tanh$ scalings; outputs are clipped to valid ranges.

\subsection{Training signals}
H-A, H-t (sup), H-J and AOD-Net minimise the $L_1$ distance between the pipeline output and the clear image; for H-A
and H-t the gradient flows through the fixed classical stages downstream of the module. \textbf{H-t (prior)} uses the
same network as H-t (sup) but no ground truth:
\begin{equation}
  \mathcal{L} = \overline{\mathrm{dark}_{15}\!\bigl(J(t)\bigr)} + \lambda_a\,\|t - t_{\text{guided}}\|_1
              + \lambda_{tv}\,\mathrm{TV}(t) + \lambda_r\,\mathcal{R}(J),
  \label{eq:prior}
\end{equation}
with $\lambda_a=1$, $\lambda_{tv}=0.1$, $\lambda_r=1$ and $\mathcal{R}$ penalising values outside $[0,1]$. Because
no ground truth is needed it trains on hazy images only, including \hybNRealHazy{} real hazy images of the tuning
split that the supervised arms never see.

\subsection{Arms}
Ten arms are evaluated: the hazy input and CLAHE (model-free controls); the DCP baseline and the tuned DCP; H-A,
H-t (sup), H-t (prior) trained with~\eqref{eq:prior}, H-J; \textbf{H-all} (H-A, H-t (sup) and H-J composed, each trained separately); and AOD-Net
(\hybAodParams{} parameters) trained here with the same data and budget.

%=========================== IV =========================================
\section{Experimental Protocol}
\textbf{Evaluation data} are the \nImages{} paired images (\dataSource) and the scene-grouped tuning/test split of the
companion study (fingerprint \texttt{\splitFingerprint}); the notebook refuses to run on a different split. All reported
numbers are on the \hybNTest{} test images. \textbf{Training data:} \hybTrainSource. Learned arms train for
\hybSteps{} steps (batch 8, $224^2$ crops, Adam $10^{-3}$ with cosine decay) with \hybSeeds{} seeds on a \hybDevice;
per-image scores are averaged over seeds before the paired statistics and the seed spread is reported.
\textbf{Metrics and statistics} follow the companion study: PSNR, SSIM~\cite{wang2004}, CIEDE2000~\cite{sharma2005},
Hautière's $e$ and $\bar r$~\cite{hautiere2008}; paired bootstrap 95\,\% CIs and Wilcoxon tests against the tuned DCP;
density terciles from the mean dark channel of the input; bright-region masks from the \emph{clear} image
(HSV $V>0.75$, $S<0.25$). The \emph{sim-to-real drop} of an arm is its mean gain on the synthetic test sets minus its
mean gain on the real ones.

%=========================== V ==========================================
\section{Results}
\subsection{Learning gain per stage (H4)}
Table~\ref{tab:main} and Fig.~\ref{fig:gain} give every arm. Relative to the tuned DCP (\hybTunedPSNR{}\,dB), the
single-stage hybrids change PSNR by \hybHADelta{} (H-A; CI [\hybHALo, \hybHAHi], $p=\hybHAP$), \hybHTsupDelta{}
(H-t sup; [\hybHTsupLo, \hybHTsupHi], $p=\hybHTsupP$) and \hybHJDelta{}\,dB (H-J; [\hybHJLo, \hybHJHi],
$p=\hybHJP$). The largest single-stage gain is at \hybBestStage{}. Composing the three modules gives
\hybHAllDelta{}\,dB ([\hybHAllLo, \hybHAllHi]) and the end-to-end AOD-Net \hybAodDelta{}\,dB ([\hybAodLo, \hybAodHi]).

\begin{table}[t]
\caption{Every arm on the test split; paired $\Delta$ against the tuned DCP; seed spread of the mean PSNR.}
\label{tab:main}
\centering\footnotesize
\includegraphics[width=\columnwidth]{tab_hybrid_main.png}
\end{table}

\begin{figure}[t]
\centering
\includegraphics[width=0.95\columnwidth]{fig_hybrid_gain_per_stage.pdf}
\caption{Learning gain per replaced stage, paired 95\,\% CIs against the tuned DCP.}
\label{fig:gain}
\end{figure}

\subsection{Density and sim-to-real transfer (H6)}
Fig.~\ref{fig:density} breaks the gains down by density stratum (Table~\ref{tab:stratum}) and by synthetic versus real
test sets. The hybrid composition gains \hybHAllSynth{}\,dB on synthetic and \hybHAllReal{}\,dB on real haze (drop
\hybHAllDrop{}); AOD-Net gains \hybAodSynth{} and \hybAodReal{}\,dB (drop \hybAodDrop{}). Per stage: H-A
\hybHADrop{}, H-t (sup) \hybHTsupDrop{}, H-J \hybHJDrop{}\,dB.

\begin{figure}[t]
\centering
\includegraphics[width=\columnwidth]{fig_hybrid_by_density.pdf}
\caption{Gain over the tuned DCP by density stratum (left) and on synthetic versus real test sets (right), paired 95\,\%
CIs.}
\label{fig:density}
\end{figure}

\begin{table}[t]
\caption{$\Delta$PSNR (dB) against the tuned DCP by density stratum.}
\label{tab:stratum}
\centering\footnotesize
\includegraphics[width=0.8\columnwidth]{tab_hybrid_by_stratum.png}
\end{table}

\subsection{The prior as a loss inside bright regions (H5)}
On the \hybBrightN{} test images whose clear image contains a bright region, Table~\ref{tab:bright} and
Fig.~\ref{fig:bright} report the per-pixel error inside the mask and elsewhere. The tuned DCP's bright/elsewhere error
ratio is \hybTunedRatio{}. The supervised refiner changes the bright-region error by \hybHTsupBDelta{}
([\hybHTsupBLo, \hybHTsupBHi], $p=\hybHTsupBP$); the prior-trained refiner by \hybHTpriorBDelta{}
([\hybHTpriorBLo, \hybHTpriorBHi], $p=\hybHTpriorBP$). Fig.~\ref{fig:tmaps} shows the four transmission estimates on
one such image.

\begin{table}[t]
\caption{Error inside bright regions vs.\ elsewhere per arm; $\Delta<0$ is better than the tuned DCP.}
\label{tab:bright}
\centering\footnotesize
\includegraphics[width=\columnwidth]{tab_hybrid_bright.png}
\end{table}

\begin{figure}[t]
\centering
\includegraphics[width=\columnwidth]{fig_hybrid_bright.pdf}
\caption{Bright-region failure per arm.}
\label{fig:bright}
\end{figure}

\begin{figure*}[t]
\centering
\includegraphics[width=0.95\textwidth]{fig_hybrid_tmaps.pdf}
\caption{Transmission on a bright-region test image: raw, guided filter, learned with ground truth, learned with the
prior as its loss; outputs through the identical recovery stage.}
\label{fig:tmaps}
\end{figure*}

\subsection{Cost}
Table~\ref{tab:cost} lists parameters, inference time and training time per arm; single-stage hybrids add
\hybHTsupParams{}--\hybHAParams{} parameters and \hybHTsupMs{}\,ms per image to the classical pipeline
(\hybTunedMs{}\,ms), against \hybAodMs{}\,ms for AOD-Net. Fig.~\ref{fig:qual} shows every arm on the same images.

\begin{table}[t]
\caption{Parameters, inference time and training time per arm.}
\label{tab:cost}
\centering\footnotesize
\includegraphics[width=0.85\columnwidth]{tab_hybrid_cost.png}
\end{table}

\begin{figure*}[t]
\centering
\includegraphics[width=0.98\textwidth]{fig_hybrid_qualitative.pdf}
\caption{Same test images, every arm (rows: worst to best for the tuned DCP).}
\label{fig:qual}
\end{figure*}

%=========================== VI =========================================
\section{Discussion}
%% AUTHORS — write from the tables. Prompts:
%%  (H4) Which stage is the bottleneck? Does the composed hybrid exceed the sum of its parts, or do the modules
%%       interfere? Is the end-to-end model ahead only because it also fixes A, t and J at once?
%%  (H5) Did the prior-as-loss refiner lower t in bright regions (over-darkening)? Compare the two refiners' bright
%%       vs elsewhere deltas — same network, same inputs, only the loss differs. What does that say about using
%%       classical priors as training signals?
%%  (H6) Which arm loses least on real haze? Is the physics scaffold acting as a regulariser? Look at Dense-Haze
%%       specifically.
%%  (cost) Where is the best accuracy per parameter / per millisecond?
\emph{To be written by the authors.}

\section{Threats to Validity}
\textbf{Budget.} \hybSteps{} steps and 17--23\,k-parameter modules make this a study of \emph{where} learning enters,
not of attainable accuracy; larger networks or longer training could reorder the arms. \textbf{Train/test mismatch.}
Training uses PyTorch classical stages with replicate padding; inference uses the NumPy stages of the companion
study. Modules are residual on the classical output, so the mismatch is a small input perturbation, but it is a
mismatch. \textbf{Data.} The supervised arms train on synthetic RESIDE haze; only H-t (prior) sees real haze, which
confounds the H5 and H6 comparisons in its favour on real sets and is why H5 is reported on all test images.
\textbf{Selection.} The scaffold's hyper-parameters were tuned on the tuning split in the companion study; no
selection touches the test split. \textbf{Seeds.} \hybSeeds{} seeds bound optimisation noise; more would tighten the
CIs.

\section{Reproducibility}
NB 4 (\texttt{dehaze-hybrid-stages}) runs on Kaggle from the attached benchmarks and NB 2's output, checks the split
fingerprint, trains every arm, writes per-image scores (\texttt{dehaze\_hybrid\_per\_image.csv}), weights per seed, and
the macros that fill this document. Seed 42 for data; seeds 0--2 for training; library versions are printed.

%=========================== VII ========================================
\section{Conclusion}
%% AUTHORS — revise once the numbers are in.
Keeping the dark channel prior pipeline as a fixed scaffold and learning one stage at a time under an identical budget,
the stage that benefits most from learning is \hybBestStage{} (\hybBestStageGain{}\,dB), the composed hybrid reaches
\hybHAllDelta{}\,dB and the end-to-end model \hybAodDelta{}\,dB over the tuned classical pipeline. A refiner trained
with the prior as its loss changes bright-region error by \hybHTpriorBDelta{} against \hybHTsupBDelta{} for its
supervised twin, and the hybrid's sim-to-real drop is \hybHAllDrop{}\,dB against \hybAodDrop{}\,dB.

\begin{thebibliography}{99}
\bibitem{narasimhan2002} S.~G.~Narasimhan and S.~K.~Nayar, ``Vision and the atmosphere,'' \emph{Int.\ J.\ Comput.\ Vis.}, vol.~48, no.~3, pp.~233--254, 2002.
\bibitem{he2009} K.~He, J.~Sun, and X.~Tang, ``Single image haze removal using dark channel prior,'' in \emph{Proc.\ CVPR}, 2009, pp.~1956--1963.
\bibitem{he2011} K.~He, J.~Sun, and X.~Tang, ``Single image haze removal using dark channel prior,'' \emph{IEEE Trans.\ Pattern Anal.\ Mach.\ Intell.}, vol.~33, no.~12, pp.~2341--2353, 2011.
\bibitem{he2010} K.~He, J.~Sun, and X.~Tang, ``Guided image filtering,'' in \emph{Proc.\ ECCV}, 2010; \emph{IEEE TPAMI}, vol.~35, no.~6, 2013.
\bibitem{wu2018} H.~Wu, S.~Zheng, J.~Zhang, and K.~Huang, ``Fast end-to-end trainable guided filter,'' in \emph{Proc.\ CVPR}, 2018.
\bibitem{cai2016} B.~Cai, X.~Xu, K.~Jia, C.~Qing, and D.~Tao, ``DehazeNet: An end-to-end system for single image haze removal,'' \emph{IEEE Trans.\ Image Process.}, vol.~25, no.~11, pp.~5187--5198, 2016.
\bibitem{ren2016} W.~Ren, S.~Liu, H.~Zhang, J.~Pan, X.~Cao, and M.-H.~Yang, ``Single image dehazing via multi-scale convolutional neural networks,'' in \emph{Proc.\ ECCV}, 2016.
\bibitem{li2017} B.~Li, X.~Peng, Z.~Wang, J.~Xu, and D.~Feng, ``AOD-Net: All-in-one dehazing network,'' in \emph{Proc.\ ICCV}, 2017, pp.~4770--4778.
\bibitem{qin2020} X.~Qin, Z.~Wang, Y.~Bai, X.~Xie, and H.~Jia, ``FFA-Net: Feature fusion attention network for single image dehazing,'' in \emph{Proc.\ AAAI}, 2020.
\bibitem{song2023} Y.~Song, Z.~He, H.~Qian, and X.~Du, ``Vision transformers for single image dehazing,'' \emph{IEEE Trans.\ Image Process.}, vol.~32, pp.~1927--1941, 2023.
\bibitem{golts2020} A.~Golts, D.~Freedman, and M.~Elad, ``Unsupervised single image dehazing using dark channel prior loss,'' \emph{IEEE Trans.\ Image Process.}, vol.~29, pp.~2692--2701, 2020.
\bibitem{li2019} B.~Li \emph{et al.}, ``Benchmarking single-image dehazing and beyond,'' \emph{IEEE Trans.\ Image Process.}, vol.~28, no.~1, pp.~492--505, 2019.
\bibitem{ancuti2018o} C.~O.~Ancuti, C.~Ancuti, R.~Timofte, and C.~De~Vleeschouwer, ``O-HAZE: A dehazing benchmark with real hazy and haze-free outdoor images,'' in \emph{Proc.\ CVPR Workshops (NTIRE)}, 2018.
\bibitem{ancuti2018i} C.~Ancuti, C.~O.~Ancuti, R.~Timofte, and C.~De~Vleeschouwer, ``I-HAZE: A dehazing benchmark with real hazy and haze-free indoor images,'' in \emph{Proc.\ ACIVS}, 2018.
\bibitem{ancuti2019} C.~O.~Ancuti, C.~Ancuti, M.~Sbert, and R.~Timofte, ``Dense-Haze: A benchmark for image dehazing with dense-haze and haze-free images,'' in \emph{Proc.\ ICIP}, 2019.
\bibitem{ancuti2020} C.~O.~Ancuti, C.~Ancuti, and R.~Timofte, ``NH-HAZE: An image dehazing benchmark with non-homogeneous hazy and haze-free images,'' in \emph{Proc.\ CVPR Workshops (NTIRE)}, 2020.
\bibitem{wang2004} Z.~Wang, A.~C.~Bovik, H.~R.~Sheikh, and E.~P.~Simoncelli, ``Image quality assessment: From error visibility to structural similarity,'' \emph{IEEE Trans.\ Image Process.}, vol.~13, no.~4, pp.~600--612, 2004.
\bibitem{sharma2005} G.~Sharma, W.~Wu, and E.~N.~Dalal, ``The CIEDE2000 color-difference formula: Implementation notes, supplementary test data, and mathematical observations,'' \emph{Color Res.\ Appl.}, vol.~30, no.~1, pp.~21--30, 2005.
\bibitem{hautiere2008} N.~Hauti\`ere, J.-P.~Tarel, D.~Aubert, and \'E.~Dumont, ``Blind contrast enhancement assessment by gradient ratioing at visible edges,'' \emph{Image Anal.\ Stereol.}, vol.~27, no.~2, pp.~87--95, 2008.
\end{thebibliography}

\end{document}
